\subsection{Manual de instalación}\label{apx:manual-instalacion}
\justifying

\subsubsection{Introducción}
Este manual describe cómo instalar y ejecutar el \emph{Sistema para la gestión de formularios sobre los cursos y autoevaluaciones docentes} de la División de Ciencias y Artes para el Diseño (CyAD) en un entorno de desarrollo local. La instalación en producción sigue los mismos pasos, cambiando únicamente el archivo \texttt{.env} y el modo de servido (\emph{gunicorn} + \emph{Nginx}, ver Subsubsección~\ref{apx:manual-inst:despliegue-produccion}).

\subsubsection{Prerrequisitos}

\begin{longtable}{p{0.30\textwidth} p{0.20\textwidth} p{0.42\textwidth}}
    \toprule
    \textbf{Software} & \textbf{Versión mínima} & \textbf{Uso} \\
    \midrule
    \endhead
    Python & 3.11+ (probado con 3.13) & Ejecución del backend Django. \\
    \midrule
    Node.js & 20+ (probado con 22 LTS) & Ejecución del frontend Vite + React. \\
    \midrule
    PostgreSQL & 14+ (probado con 17) & Base de datos relacional. \\
    \midrule
    Git & 2.40+ & Clonado del repositorio. \\
    \midrule
    \emph{npm} (incluido con Node) & 10+ & Gestor de paquetes JS. \\
    \bottomrule
\end{longtable}

En Windows se recomienda utilizar \emph{Git Bash} o \emph{PowerShell}. En Linux/macOS los comandos son idénticos, cambiando la ruta del entorno virtual (\texttt{bin/} en lugar de \texttt{Scripts/}).

\textbf{Nota importante para Windows:} el instalador oficial de PostgreSQL no agrega \texttt{psql} al \texttt{PATH} del sistema. Si al escribir \texttt{psql --version} obtiene ``command not found'', agregue manualmente la carpeta \texttt{C:\textbackslash Program Files\textbackslash PostgreSQL\textbackslash 17\textbackslash bin} (ajuste al número de versión que tenga) a la variable de entorno \texttt{Path} de Windows, o invoque \texttt{psql} por ruta completa.

\paragraph{Requisitos mínimos de hardware}\mbox{}\\
Las siguientes cifras se derivan de mediciones sobre un entorno con \texttt{seed\_demo} cargado, ejecutando \texttt{backend/scripts/medir\_hardware.py} (descrito al final de este apartado): el proceso de Django consume aproximadamente 120~MB de RAM en reposo, PostgreSQL alrededor de 80~MB con la carga de demostración, y el \emph{build} del frontend (\texttt{npm run build}) alcanza un pico de \textasciitilde 250~MB durante la compilación (no en tiempo de ejecución).

\begin{longtable}{p{0.30\textwidth} p{0.20\textwidth} p{0.20\textwidth} p{0.22\textwidth}}
    \toprule
    \textbf{Escenario} & \textbf{RAM mínima} & \textbf{RAM recomendada} & \textbf{Disco} \\
    \midrule
    \endhead
    Desarrollo local (una sola persona) & 4~GB & 8~GB & 5~GB libres \\
    \midrule
    Servidor pequeño (\textless\,50 profesores, un solo periodo activo) & 2~GB & 4~GB & 10~GB libres \\
    \midrule
    Servidor CyAD completo (\textgreater\,300 profesores, histórico multi-periodo) & 4~GB & 8~GB & 20~GB libres \\
    \bottomrule
\end{longtable}

Sobre CPU: cualquier procesador x86-64 de doble núcleo publicado en la última década es suficiente para el runtime; el cuello de botella es la memoria durante \texttt{npm run build}, que sólo ocurre al desplegar una nueva versión, no en operación diaria. En un servidor dedicado bastan \emph{dos vCPU}; para desarrollo local la máquina que edita el código suele tener recursos de sobra.

Ancho de banda: al ser una aplicación de tipo formulario, el consumo por sesión es del orden de decenas de KB por interacción. Cualquier enlace institucional (\textgreater\,10~Mbps) resulta holgado.

\paragraph{Reproducibilidad de las cifras}\mbox{}\\
Las tres cifras anteriores (Django en reposo, PostgreSQL con demo, pico de \texttt{npm run build}) se pueden reproducir sobre cualquier host ejecutando el \emph{script} versionado \texttt{backend/scripts/medir\_hardware.py}. El \emph{script}:

\begin{itemize}
    \item arranca \texttt{manage.py runserver --noreload} en un subproceso, sirve una petición a \texttt{/api/schema/} para forzar la inicialización de la app y muestrea el \texttt{RSS} del proceso con \texttt{psutil};
    \item enumera todos los procesos \texttt{postgres} visibles al usuario y suma su \texttt{RSS} (refleja la carga vigente de la base; ideal medirlo después de correr \texttt{python scripts/seed\_demo.py});
    \item ejecuta \texttt{npm run build} en \texttt{frontend/} y muestrea el RSS del árbol de procesos cada 200~ms hasta que termina, guardando el máximo.
\end{itemize}

Al finalizar imprime una tabla legible en la consola y deja el detalle en \texttt{backend/scripts/medir\_hardware.out.json}, con la marca de tiempo y el hardware del host (\texttt{platform.uname()}, RAM total, vCPU). Requiere \texttt{psutil} (incluido en \texttt{backend/requirements.txt}) y que \texttt{frontend/node\_modules/} exista (correr \texttt{npm install} antes).

\subsubsection{Clonado del repositorio}

\begin{quote}
\ttfamily
git clone --depth 1 \textless URL\_del\_repositorio\textgreater\ proyecto\_terminal\_cyad \\
cd proyecto\_terminal\_cyad
\end{quote}

\subsubsection{Instalación mínima sin datos de demo}
\label{apx:manual-inst:install-limpia}

Esta es la ruta recomendada para un despliegue institucional real. Deja el sistema con los catálogos oficiales de CyAD (departamentos, licenciaturas, posgrados, áreas y UEAs) y una única cuenta de administrador; no crea profesores de prueba, formularios ficticios ni respuestas simuladas. Es la única ruta apropiada para producción; el flujo con datos de demostración se documenta después (Subsubsección~\ref{apx:manual-inst:datos-demo}) y no debe ejecutarse en un ambiente productivo.

\paragraph{1. Entorno virtual y dependencias del backend}\mbox{}\\
\begin{quote}
\ttfamily
cd backend \\
python -m venv .venv \\
.venv\textbackslash Scripts\textbackslash activate\ \ \# Windows \\
source .venv/bin/activate\ \ \# Linux/macOS \\
pip install -r requirements.txt
\end{quote}

\paragraph{2. Archivo \texttt{.env} del backend}\mbox{}\\
\textbf{Ubicación:} el archivo \texttt{.env} que Django lee debe estar en la carpeta \texttt{backend/} (donde vive \texttt{manage.py}), \emph{no} en la raíz del repositorio. El repositorio incluye \texttt{backend/.env.example} como plantilla; cópielo y ajuste los valores:

\begin{quote}
\ttfamily
cp .env.example .env\ \ \# desde backend/
\end{quote}

Contenido esperado (valores por defecto de la plantilla; ajuste \texttt{DB\_PASSWORD} y \texttt{SECRET\_KEY} a su entorno):

\begin{Verbatim}[fontsize=\small, frame=single]
DEBUG=True
SECRET_KEY=cambia-esta-clave-por-una-aleatoria-larga
ALLOWED_HOSTS=localhost,127.0.0.1
DB_NAME=proyecto_terminal_cyad
DB_USER=postgres
DB_PASSWORD=tu_password_de_postgres
DB_HOST=localhost
DB_PORT=5432
JWT_ACCESS_MINUTES=60
JWT_REFRESH_DAYS=7
CORS_ALLOWED_ORIGINS=http://localhost:5173,http://127.0.0.1:5173
\end{Verbatim}

Para producción se debe establecer \texttt{DEBUG=False}, definir \texttt{ALLOWED\_HOSTS} con el dominio real y generar una \texttt{SECRET\_KEY} con \texttt{python -c "from django.core.management.utils import get\_random\_secret\_key; print(get\_random\_secret\_key())"} (más detalle en Subsubsección~\ref{apx:manual-inst:despliegue-produccion}).

\paragraph{3. Crear la base de datos vacía}\mbox{}\\
El proyecto utiliza por defecto el usuario \texttt{postgres} y una base de datos llamada \texttt{proyecto\_terminal\_cyad} (ver \texttt{backend/.env.example}). La forma recomendada es dejar que el script auxiliar cree la base automáticamente:

\begin{quote}
\ttfamily
python scripts/create\_db.py
\end{quote}

Si prefiere crearla manualmente con \texttt{psql}:

\begin{quote}
\ttfamily
psql -U postgres \\
CREATE DATABASE proyecto\_terminal\_cyad;
\end{quote}

También puede definir un usuario dedicado con permisos limitados (recomendado en despliegues compartidos) y ajustar \texttt{DB\_USER}/\texttt{DB\_PASSWORD} en \texttt{.env} en consecuencia:

\begin{quote}
\ttfamily
CREATE USER cyad\_user WITH PASSWORD 'cyad\_pass'; \\
GRANT ALL PRIVILEGES ON DATABASE proyecto\_terminal\_cyad TO cyad\_user;
\end{quote}

\paragraph{4. Aplicar migraciones}\mbox{}\\
\begin{quote}
\ttfamily
python manage.py migrate
\end{quote}

\paragraph{5. Crear la primera cuenta de administrador}\mbox{}\\
Con el comando interactivo de Django:

\begin{quote}
\ttfamily
python manage.py createsuperuser
\end{quote}

Se pedirán correo, nombre y contraseña. Este usuario tendrá rol \texttt{ADMIN} y podrá crear el resto de cuentas desde la interfaz web.

\paragraph{6. Cargar los catálogos base (fixtures JSON versionados)}\mbox{}\\
Departamentos, licenciaturas, posgrados y áreas están precargados en fixtures JSON dentro del repositorio, y se aplican con el comando \texttt{loaddata} de Django. Ejecutar en el orden que dictan las claves foráneas:

\begin{quote}
\ttfamily
python manage.py loaddata catalogos/fixtures/departamentos.json \\
python manage.py loaddata catalogos/fixtures/licenciaturas.json \\
python manage.py loaddata catalogos/fixtures/posgrados.json \\
python manage.py loaddata catalogos/fixtures/areas.json
\end{quote}

Los fixtures son idempotentes (usan clave primaria explícita): se pueden re-ejecutar sin duplicar registros.

\paragraph{7. Cargar las UEAs oficiales}\mbox{}\\
Los CSV oficiales de coordinación académica (\texttt{ueas\_licenciatura.csv} con las UEAs de DCG, ARQ, DI y DiPS; \texttt{ueas\_posgrado.csv} con las UEAs de PPCDA, PDB y PDEU) viven versionados en \texttt{backend/catalogos/fixtures/}. Se cargan con el comando de gestión \texttt{cargar\_catalogos\_csv}:

\begin{quote}
\ttfamily
python manage.py cargar\_catalogos\_csv \\
\ \ --ueas-csv catalogos/fixtures/ueas\_licenciatura.csv \\
python manage.py cargar\_catalogos\_csv \\
\ \ --ueas-csv catalogos/fixtures/ueas\_posgrado.csv
\end{quote}

El comando valida claves duplicadas, reporta filas rechazadas y resuelve el área de cada UEA contra los pares \texttt{(nombre, descripción)} de \texttt{areas.json}; las áreas que no encuentre las lista como \emph{warning} y la UEA queda con \texttt{area = NULL} (nunca se crea un área nueva de forma silenciosa). La carga es idempotente por el \texttt{update\_or\_create} interno.

Si coordinación entrega en el futuro un CSV incremental (por ejemplo cuando cambie un plan de estudios), se carga con el mismo comando apuntando \texttt{--ueas-csv} a la nueva ruta; el efecto es un \emph{upsert} por \texttt{(clave, programa)}.

\paragraph{8. Crear el periodo académico vigente}\mbox{}\\
Los periodos (25-P, 26-I, …) \emph{no} se precargan como fixture porque son datos temporales que envejecen mal (un fixture con periodos históricos quedaría obsoleto tras unos meses). El primer administrador los captura desde la interfaz una vez que el sistema está corriendo, en \texttt{/admin/catalogos/periodos}: se ingresa clave, fecha de inicio, fecha de fin y se activan los \emph{flags} por recurso (\texttt{activo\_cartas}, \texttt{activo\_requisitos}, \texttt{activo\_autoevaluacion}) según la ventana de captura vigente. Ver §~``Módulo 3: habilitación de periodos por recurso'' del capítulo de Desarrollo (\S~\ref{sec:desarrollo}) para el detalle de estos \emph{flags}.

\paragraph{9. Arrancar el servidor de desarrollo del backend}\mbox{}\\
\begin{quote}
\ttfamily
python manage.py runserver 8000
\end{quote}

El backend queda disponible en \texttt{http://localhost:8000}. La documentación interactiva de la API está en \texttt{http://localhost:8000/api/docs/}.

\paragraph{10. Crear el formulario de autoevaluación institucional}\mbox{}\\
Con el backend corriendo, iniciar sesión con el superuser del paso 5 y, desde \texttt{Autoevaluación $\to$ Nuevo formulario}, construir el formulario del periodo vigente (secciones, preguntas, opciones puntuables y niveles de desempeño). Publicarlo cuando la suma de pesos de las secciones alcance 100~\%.

Los usuarios de profesor se crean posteriormente desde \texttt{Profesores $\to$ Nuevo profesor} en la interfaz de administrador; cada creación genera un correo + contraseña temporal que el administrador entrega personalmente al docente.

\subsubsection{Datos de demostración (opcional)}
\label{apx:manual-inst:datos-demo}

\textbf{Advertencia:} este bloque es exclusivo del \emph{desarrollo local} y de la evaluación en escritorio; \emph{no debe ejecutarse en un ambiente productivo}. Crea un profesor de prueba, un formulario ya publicado y respuestas ficticias que después será tedioso eliminar.

\paragraph{Bootstrap unificado (Windows PowerShell)}\mbox{}\\
Como atajo, el repositorio incluye \texttt{scripts/bootstrap.ps1} que orquesta \emph{creación de BD}, \emph{migraciones} y \emph{siembra de datos de demo} en un único comando. Debe ejecutarse desde la raíz del repositorio, con el \texttt{.venv} ya activo y \texttt{backend/.env} configurado:

\begin{quote}
\ttfamily
.\textbackslash scripts\textbackslash bootstrap.ps1
\end{quote}

El script imprime tres pasos numerados (\texttt{[1/3]} crear BD, \texttt{[2/3]} migrar, \texttt{[3/3]} seed) y al finalizar recuerda cómo arrancar el backend y el frontend en dos terminales distintas. Falla temprano si falta \texttt{backend/.env} o si el árbol del repositorio no está completo.

\paragraph{Seed manual (Linux, macOS, o alternativa a Bootstrap)}\mbox{}\\
Con la base de datos ya creada y las migraciones aplicadas (pasos 3 y 4 de la Subsubsección~\ref{apx:manual-inst:install-limpia}), ejecute:

\begin{quote}
\ttfamily
python scripts/seed\_demo.py
\end{quote}

El \emph{seed} es idempotente y crea catálogos, UEAs, usuarios de prueba, periodo activo (\texttt{26-I}) y el formulario de autoevaluación en un solo comando. Al finalizar, existirán las siguientes cuentas:

\begin{longtable}{p{0.30\textwidth} p{0.35\textwidth} p{0.30\textwidth}}
    \toprule
    \textbf{Rol} & \textbf{Correo} & \textbf{Contraseña} \\
    \midrule
    \endhead
    Administrador & admin@cyad.uam.mx & Admin1234! \\
    \midrule
    Profesor & profesor@cyad.uam.mx & Profesor1234! \\
    \midrule
    Profesor (segundo) & profesor2@cyad.uam.mx & Profesor1234! \\
    \bottomrule
\end{longtable}

Estas contraseñas deben cambiarse antes del despliegue a producción (y, de hecho, la base entera debería recrearse limpia siguiendo la Subsubsección~\ref{apx:manual-inst:install-limpia}).

\subsubsection{Configuración del frontend}

En una terminal aparte:

\begin{quote}
\ttfamily
cd frontend \\
npm install \\
npm run dev
\end{quote}

El frontend queda disponible en \texttt{http://localhost:5173}. Vite realiza \emph{hot reload} automático al editar archivos fuente.

Si el backend no está en \texttt{localhost:8000}, cree \texttt{frontend/.env.local} con:

\begin{Verbatim}[fontsize=\small, frame=single]
VITE_API_URL=http://tu-servidor:8000/api/v1
\end{Verbatim}

\textbf{Nota:} el nombre correcto de la variable es \texttt{VITE\_API\_URL} (verificable en \texttt{frontend/src/api/client.js}).

\subsubsection{Verificación}

\begin{enumerate}
    \item Abra \texttt{http://localhost:5173/} — debe mostrarse la vista pública ``Explorar''.
    \item Inicie sesión con el superuser creado en el paso 5 de la Subsubsección~\ref{apx:manual-inst:install-limpia} (o, si sembró la demo, con \texttt{admin@cyad.uam.mx / Admin1234!}). Debe redirigir al \emph{dashboard} de administrador.
    \item Si sembró la demo, cierre sesión e inicie con \texttt{profesor@cyad.uam.mx / Profesor1234!}. Debe redirigir al \emph{dashboard} de profesor.
    \item Ejecute la suite de pruebas: \texttt{pytest} desde \texttt{backend/}. Todas deben pasar.
    \item \emph{Opcional}: con el backend corriendo y el seed cargado, ejecute \texttt{python scripts/smoke\_test.py} desde \texttt{backend/}. Verifica de extremo a extremo, vía HTTP, el \emph{login} de \texttt{admin} y \texttt{profesor}, \texttt{/auth/me}, catálogos, formularios disponibles, creación de respuesta, creación de carta temática y los reportes de administrador.
\end{enumerate}

\subsubsection{Comandos administrativos útiles}

\begin{longtable}{p{0.55\textwidth} p{0.38\textwidth}}
    \toprule
    \textbf{Comando} & \textbf{Uso} \\
    \midrule
    \endhead
    \texttt{python manage.py cargar\_catalogos\_csv --ueas-csv <archivo>} & Importar UEAs adicionales desde CSV. \\
    \midrule
    \texttt{python manage.py createsuperuser} & Crear cuenta administradora adicional. \\
    \midrule
    \texttt{python manage.py dumpdata <app>} & Exportar datos a JSON. \\
    \midrule
    \texttt{python scripts/medir\_hardware.py} & Medir RAM de Django, PostgreSQL y pico de \texttt{npm run build}. \\
    \midrule
    \texttt{pytest --cov=. --cov-report=html} & Ejecutar pruebas con cobertura. \\
    \midrule
    \texttt{npm run build} (frontend) & Compilar producción estática (\texttt{dist/}). \\
    \bottomrule
\end{longtable}

\subsubsection{Despliegue en servidor de producción}
\label{apx:manual-inst:despliegue-produccion}

Esta subsubsección describe el despliegue en un servidor Linux (Ubuntu 22.04 LTS o Debian 12) que atienda las peticiones institucionales. El manual asume que ya se cumplen los prerrequisitos (Subsubsección~\ref{apx:manual-inst:install-limpia} inclusive: base limpia, superuser creado, catálogos cargados).

\paragraph{Topología recomendada}\mbox{}\\
\begin{Verbatim}[fontsize=\small, frame=single]
        Internet (HTTPS 443)
                |
                v
     +--------------------+
     |       Nginx        |    Sirve /  -> frontend/dist  (SPA estático)
     | (proxy inverso +   |    /api/    -> gunicorn (socket UNIX)
     |  archivos estát.)  |    /admin/  -> gunicorn (socket UNIX)
     +---------+----------+    /static/ -> backend/staticfiles
               |
               v (unix:/run/cyad.sock)
     +--------------------+
     |  gunicorn (3 wk)   |    Django / DRF
     |  config.wsgi:app   |
     +---------+----------+
               |
               v (TCP 5432)
     +--------------------+
     |    PostgreSQL 17   |    Usuario dedicado cyad_prod
     +--------------------+
\end{Verbatim}

Todos los servicios corren en la misma máquina en el caso más simple; se pueden separar la BD y el \emph{app server} sin cambios en el código, ajustando \texttt{DB\_HOST} en el \texttt{.env}.

\paragraph{Variables de entorno de producción}\mbox{}\\
El archivo \texttt{backend/.env} para producción difiere del de desarrollo en cinco campos críticos:

\begin{Verbatim}[fontsize=\small, frame=single]
DEBUG=False
SECRET_KEY=<clave-larga-generada-abajo>
ALLOWED_HOSTS=cyad.uam.mx
CORS_ALLOWED_ORIGINS=https://cyad.uam.mx
DB_NAME=cyad_prod
DB_USER=cyad_prod
DB_PASSWORD=<password-fuerte>
DB_HOST=localhost
DB_PORT=5432
JWT_ACCESS_MINUTES=15
JWT_REFRESH_DAYS=1
\end{Verbatim}

Generación de \texttt{SECRET\_KEY}:

\begin{quote}\ttfamily
python -c "from django.core.management.utils import get\_random\_secret\_key; \\
print(get\_random\_secret\_key())"
\end{quote}

La ventana de \texttt{JWT\_ACCESS\_MINUTES=15} y \texttt{JWT\_REFRESH\_DAYS=1} reduce el impacto de un token filtrado; el frontend refresca el \emph{access token} automáticamente en segundo plano mientras el \emph{refresh} siga vigente.

El archivo \texttt{.env} debe tener permisos restrictivos:

\begin{quote}\ttfamily
chmod 600 backend/.env \\
chown www-data:www-data backend/.env
\end{quote}

\paragraph{Base de datos en producción}\mbox{}\\
Crear un rol \emph{no-superuser} dedicado y una base propiedad de ese rol:

\begin{quote}\ttfamily
sudo -u postgres psql \\
CREATE USER cyad\_prod WITH PASSWORD '\textless password-fuerte\textgreater'; \\
CREATE DATABASE cyad\_prod OWNER cyad\_prod; \\
\textbackslash q
\end{quote}

Endurecer \texttt{pg\_hba.conf} para permitir sólo conexiones locales autenticadas por contraseña al usuario \texttt{cyad\_prod} (no permita accesos remotos al usuario \texttt{postgres}).

Respaldo diario recomendado vía \emph{cron}:

\begin{Verbatim}[fontsize=\small, frame=single]
# /etc/cron.d/cyad-backup
0 3 * * * postgres pg_dump -Fc cyad_prod \
  > /var/backups/cyad/cyad_$(date +\%Y\%m\%d).dump
\end{Verbatim}

Rotar respaldos con \texttt{logrotate} o un simple \texttt{find /var/backups/cyad -mtime +30 -delete}.

\paragraph{Backend con gunicorn + systemd}\mbox{}\\
Instalar gunicorn en el mismo entorno virtual:

\begin{quote}\ttfamily
pip install gunicorn
\end{quote}

Crear la unidad \texttt{/etc/systemd/system/cyad-backend.service}:

\begin{Verbatim}[fontsize=\small, frame=single]
[Unit]
Description=CyAD Django backend (gunicorn)
After=network.target postgresql.service
Requires=postgresql.service

[Service]
Type=notify
User=www-data
Group=www-data
WorkingDirectory=/srv/cyad/proyecto_terminal_cyad/backend
Environment="PATH=/srv/cyad/proyecto_terminal_cyad/backend/.venv/bin"
ExecStart=/srv/cyad/proyecto_terminal_cyad/backend/.venv/bin/gunicorn \
    config.wsgi:application \
    --workers 3 \
    --bind unix:/run/cyad.sock \
    --access-logfile /var/log/cyad/access.log \
    --error-logfile /var/log/cyad/error.log
Restart=on-failure

[Install]
WantedBy=multi-user.target
\end{Verbatim}

Habilitar y arrancar:

\begin{quote}\ttfamily
sudo mkdir -p /var/log/cyad \\
sudo chown www-data:www-data /var/log/cyad \\
sudo systemctl daemon-reload \\
sudo systemctl enable --now cyad-backend \\
sudo systemctl status cyad-backend
\end{quote}

\paragraph{Archivos estáticos del backend}\mbox{}\\
El Django admin y \texttt{drf-spectacular} sirven CSS/JS estáticos que Nginx debe entregar directamente. Recolectarlos en \texttt{backend/staticfiles/} (ruta ya definida como \texttt{STATIC\_ROOT} en \texttt{config/settings.py}):

\begin{quote}\ttfamily
python manage.py collectstatic --noinput
\end{quote}

Debe re-ejecutarse en cada actualización que añada estáticos nuevos.

\paragraph{Frontend en producción}\mbox{}\\
Compilar los estáticos del SPA:

\begin{quote}\ttfamily
cd frontend \\
npm ci \\
npm run build
\end{quote}

\texttt{npm ci} respeta \texttt{package-lock.json} exactamente y es más rápido que \texttt{npm install} para despliegues reproducibles. El resultado queda en \texttt{frontend/dist/}, que Nginx servirá como archivos estáticos.

Antes de compilar, crear \texttt{frontend/.env.production} con la URL pública de la API:

\begin{Verbatim}[fontsize=\small, frame=single]
VITE_API_URL=https://cyad.uam.mx/api/v1
\end{Verbatim}

\paragraph{Nginx como reverse proxy}\mbox{}\\
Crear \texttt{/etc/nginx/sites-available/cyad.conf}:

\begin{Verbatim}[fontsize=\small, frame=single]
server {
    listen 80;
    server_name cyad.uam.mx;

    # SPA: sirve archivos estáticos y hace fallback a index.html
    root /srv/cyad/proyecto_terminal_cyad/frontend/dist;
    index index.html;

    location / {
        try_files $uri $uri/ /index.html;
    }

    # Backend (API + admin) via socket UNIX
    location /api/ {
        proxy_pass http://unix:/run/cyad.sock;
        proxy_set_header Host              $host;
        proxy_set_header X-Real-IP         $remote_addr;
        proxy_set_header X-Forwarded-For   $proxy_add_x_forwarded_for;
        proxy_set_header X-Forwarded-Proto $scheme;
    }

    location /admin/ {
        proxy_pass http://unix:/run/cyad.sock;
        proxy_set_header Host              $host;
        proxy_set_header X-Real-IP         $remote_addr;
        proxy_set_header X-Forwarded-For   $proxy_add_x_forwarded_for;
        proxy_set_header X-Forwarded-Proto $scheme;
    }

    # Estáticos del backend (Django admin, drf-spectacular)
    location /static/ {
        alias /srv/cyad/proyecto_terminal_cyad/backend/staticfiles/;
        expires 30d;
    }

    client_max_body_size 10M;
}
\end{Verbatim}

Habilitar y recargar:

\begin{quote}\ttfamily
sudo ln -s /etc/nginx/sites-available/cyad.conf \\
\ \ \ \ \ \ /etc/nginx/sites-enabled/ \\
sudo nginx -t \\
sudo systemctl reload nginx
\end{quote}

\paragraph{HTTPS con Let's Encrypt}\mbox{}\\
Con \emph{certbot} instalado (\texttt{sudo apt install certbot python3-certbot-nginx}):

\begin{quote}\ttfamily
sudo certbot --nginx -d cyad.uam.mx
\end{quote}

Certbot modifica el bloque de Nginx anterior para añadir escucha en \texttt{:443} con el certificado emitido y una redirección \texttt{80 $\to$ 443}. La renovación se instala automáticamente vía \emph{systemd timer}; verificar con \texttt{sudo systemctl status certbot.timer}.

\paragraph{Actualizaciones (redeploy)}\mbox{}\\
Cada nueva versión del código se despliega con:

\begin{quote}\ttfamily
cd /srv/cyad/proyecto\_terminal\_cyad \\
git pull \\
cd backend \\
source .venv/bin/activate \\
pip install -r requirements.txt \\
python manage.py migrate \\
python manage.py collectstatic --noinput \\
cd ../frontend \\
npm ci \\
npm run build \\
sudo systemctl restart cyad-backend \\
sudo systemctl reload nginx
\end{quote}

Es recomendable hacer \texttt{pg\_dump} antes de cada \texttt{migrate} en producción.

\paragraph{Nota para despliegues en Windows Server}\mbox{}\\
Si el hosting institucional obliga a usar Windows Server, gunicorn no está disponible. La alternativa recomendada es \emph{waitress}:

\begin{quote}\ttfamily
pip install waitress \\
waitress-serve --port=8000 config.wsgi:application
\end{quote}

Como \emph{reverse proxy} se usa IIS con el módulo \emph{URL Rewrite} + \emph{Application Request Routing}, o bien \emph{Nginx for Windows} con una configuración equivalente a la mostrada arriba. El servicio se registra con \texttt{sc.exe create} o mediante \emph{NSSM}. El resto de la configuración (\texttt{.env}, base de datos, \texttt{collectstatic}, \texttt{npm run build}) es idéntica.

\subsubsection{Solución de problemas}

\begin{itemize}
    \item \textbf{Error \texttt{FATAL: password authentication failed}}: verifique \texttt{DB\_USER} y \texttt{DB\_PASSWORD} en \texttt{.env} y que el usuario tenga permisos sobre la base de datos.
    \item \textbf{Error CORS al iniciar sesión desde el frontend}: verifique que \texttt{CORS\_ALLOWED\_ORIGINS} incluya \texttt{http://localhost:5173}.
    \item \textbf{\texttt{429 Too Many Requests} al iniciar sesión}: se activó el \emph{rate limiting} (5 intentos por minuto). Espere 60 segundos y reintente.
    \item \textbf{Puerto 5173 ocupado}: Vite usará automáticamente el siguiente puerto libre; verifique la URL impresa al inicio.
    \item \textbf{\texttt{cargar\_catalogos\_csv} reporta áreas no encontradas}: significa que el CSV trae un par \texttt{(area\_nombre, area\_descripcion)} que no está precargado en \texttt{areas.json}. Agregue el área faltante (ya sea editando \texttt{areas.json} y volviendo a hacer \texttt{loaddata}, o desde \texttt{/admin/catalogos/areas} en la UI) y vuelva a ejecutar el comando; las UEAs afectadas se actualizarán con el área correcta.
    \item \textbf{\texttt{psql: command not found} en Windows}: agregue \texttt{C:\textbackslash Program Files\textbackslash PostgreSQL\textbackslash<versión>\textbackslash bin} al \texttt{PATH}, o llame a \texttt{psql} por ruta completa. Alternativamente, use \emph{pgAdmin} o el script \texttt{scripts/create\_db.py}, que utiliza el driver Python y no requiere \texttt{psql} en el PATH.
    \item \textbf{\texttt{scripts/create\_db.py} responde \texttt{'utf-8' codec can't decode byte \dots: invalid continuation byte}}: es una traducción cifrada de un fallo de conexión de \emph{libpq} cuyo mensaje viene en la codificación regional de Windows (\emph{cp1252}) y no en \emph{utf-8}. La causa típica es autenticación fallida contra PostgreSQL: revise \texttt{DB\_USER} y \texttt{DB\_PASSWORD} en \texttt{backend/.env} y confirme que ese rol pueda conectarse a la base \texttt{postgres} (\texttt{psql -U <usuario> -d postgres}). El mismo síntoma aparece si el servicio de PostgreSQL no está corriendo.
\end{itemize}
